%% sec/results.tex — §5 Results: Cost-Quality Pareto
%% Source: paper/sec5_results.md
%% NB: all quantitative cells are TBD-from-D6 placeholders. mlcad_runner D6-D8 fills.

\section{Results: Cost--Quality Pareto}
\label{sec:results}

\subsection{Overall result (TL;DR)}
\label{sec:results-tldr}

Across eleven checkpoints drawn from seven vendor families on
LC\_VCO\_20G, all 11 LLMs meet the full spec contract on the
median seed (33/33 cells PASS); the cost--quality Pareto frontier
spans \textbf{2.4 orders of magnitude in USD/run} (a 270$\times$
spread) while success-rate variance collapses to 0 percentage
points. The four China-frontier vendors collectively occupy the
\textbf{low-cost-high-success} region of the Pareto, with DeepSeek
V4-flash (\$0.0028/run) the single cheapest checkpoint of the
eleven---demonstrating that domestic-China inference is not a
single anecdote but a cross-vendor pattern.

\subsection{Figure 2 --- Cost--Quality Pareto scatter}
\label{sec:results-pareto-fig}

The Pareto scatter figure (rendering deferred to D9 once D8 data
lands; raw points generated by the benchmark runner from
\filename{paper/data/extracted\_logs.csv}) specifies the scatter:

\begin{itemize}
  \item \textbf{x-axis} --- USD/run, log scale, range TBD spanning
    the cost spread of the eleven checkpoints under the
    2026-05-01 pricing snapshot.
  \item \textbf{y-axis} --- success rate (median across seeds)
    $\in[0,1]$.
  \item \textbf{Markers} --- one per (LLM, seed) run; per-LLM
    median highlighted with a larger filled marker.
  \item \textbf{Shape} encodes country: $\circ$~$=$~US,
    $\bullet$~$=$~China.
  \item \textbf{Colour} encodes vendor family: one distinct colour
    each for Anthropic, OpenAI, Google, Moonshot, MiniMax, Xiaomi,
    and DeepSeek (seven colours total). Intra-vendor checkpoints
    share colour and are distinguished by marker fill / size as
    documented in the caption.
  \item \textbf{Pareto frontier} drawn in bold over the convex-hull
    lower-right envelope (max success per cost).
  \item \textbf{Annotations}: name-labelled annotations on the
    US-frontier checkpoints (Opus 4.7, GPT-5.5, Gemini 2.x) and
    the China-frontier cohort (Kimi K2.5, MiniMax M2.7, MiMo
    v2.5-pro, DeepSeek V4-pro).
  \item \textbf{Inset} (top-left): on this workload all 33 cells
    converge to success rate $1.00$, collapsing the inset to a
    horizontal line; an inset is retained in the figure spec only
    so a future noisier workload can re-use the rendering pipeline.
\end{itemize}

\todo[inline]{Figure 2 placeholder; D6 numeric ingest complete (33/33 cells), render script not yet wired -- camera-ready hand-off.}

\subsection{Per-LLM benchmark data}
\label{sec:results-table}

\begin{table*}[t]
  \centering
  \caption{Per-LLM benchmark data on LC\_VCO\_20G. Cells are
    median across 3 seeds; per-cell IQR in supplementary appendix.
    USD/run computed against the 2026-05-01 public pricing
    snapshot (appendix).}
  \label{tab:percellm}
  \footnotesize
  \renewcommand{\arraystretch}{0.88}
  \setlength{\tabcolsep}{3pt}
  \begin{tabular}{l l l l l l l}
    \toprule
    \textbf{LLM (vendor, country)} &
      \makecell[l]{\textbf{Success rate}\\\textbf{(median, IQR)}} &
      \makecell[l]{\textbf{Iter-to-}\\\textbf{converge}} &
      \makecell[l]{\textbf{Input}\\\textbf{tokens}} &
      \makecell[l]{\textbf{Output}\\\textbf{tokens}} &
      \makecell[l]{\textbf{Reasoning}\\\textbf{tokens}} &
      \makecell[l]{\textbf{USD /}\\\textbf{run}} \\
    \midrule
    Opus 4.7 (Anthropic, US)            & 1.00 & 3 & 30.4k & 2.7k  & \textit{n/a} & \$0.66   \\
    Sonnet 4.6 (Anthropic, US)          & 1.00 & 2 & 11.7k & 1.9k  & \textit{n/a} & \$0.063  \\
    Haiku 4.5 (Anthropic, US)           & 1.00 & 3 & 25.4k & 1.9k  & \textit{n/a} & \$0.028  \\
    GPT-5.5 (OpenAI, US)                & 1.00 & 2 & 9.8k  & 1.3k  & 476          & \$0.030  \\
    GPT-5.4-mini (OpenAI, US)           & 1.00 & 2 & 9.9k  & 997   & 0            & \$0.0045 \\
    Gemini 2.x (Google, US)             & 1.00 & 2 & 12.0k & 868   & \textit{n/a} & \$0.024  \\
    Kimi K2.5 (Moonshot, China)         & 1.00 & 2 & 9.6k  & 7.2k  & 0            & \$0.024  \\
    MiniMax M2.7 (MiniMax, China)       & 1.00 & 4 & 38.1k & 5.5k  & 2.5k         & \$0.022  \\
    MiMo v2.5-pro (Xiaomi, China)       & 1.00 & 2 & 11.8k & 4.9k  & 4.1k         & \$0.019  \\
    DeepSeek V4-pro (DeepSeek, China)   & 1.00 & 2 & 9.9k  & 6.2k  & 5.2k         & \$0.018  \\
    DeepSeek V4-flash (DeepSeek, China) & 1.00 & 3 & 21.3k & 3.6k  & 3.2k         & \$0.0028 \\
    \bottomrule
  \end{tabular}
\end{table*}

\subsection{Reasoning-token premium --- SP3 sharp evidence}
\label{sec:results-premium}

A sharp single observation anticipated in this work is the
\textbf{reasoning-token premium} on vendors that bill reasoning
tokens in their per-token rate but report them in a separate
channel. Buyers estimating cost from the \emph{visible response
text}---and applying the public per-token rate---systematically
under-budget by a vendor-specific fractional ratio
$r/(1-r)$. We define
$r \coloneqq \texttt{reasoning\_tokens} / \texttt{completion\_tokens}$,
where \texttt{completion\_tokens} is the billed completion total
surfaced by \filename{\_normalize\_usage} and \textbf{includes}
reasoning.

\begin{itemize}
  \item For \textbf{Kimi K2.5}, \textbf{MiniMax M2.7},
    \textbf{DeepSeek V4-pro / V4-flash}, and \textbf{Xiaomi MiMo
    v2.5-pro} (vendors exposing \texttt{reasoning\_content} and
    tracked as a separate \texttt{reasoning\_tokens} field by
    \filename{\_normalize\_usage} at
    \filename{src/llm\_client.py:21-80}), the median reasoning
    fractions are Kimi K2.5 $r = 0.00$ (null in default mode),
    MiniMax M2.7 $r = 0.52$, MiMo v2.5-pro $r = 0.73$, DeepSeek
    V4-pro $r = 0.89$, DeepSeek V4-flash $r = 0.83$.

  \item For \textbf{OpenAI GPT-5.5 / GPT-5.4-mini} the
    \texttt{reasoning\_tokens} counter is reported as a sub-field
    of completion tokens; the premium is GPT-5.5 $r = 0.37$,
    GPT-5.4-mini $r = 0.00$ (no reasoning emitted).

  \item For \textbf{Anthropic Opus 4.7 / Sonnet 4.6 / Haiku 4.5}
    and \textbf{Google Gemini 2.x}: per-iteration reasoning-token
    accounting is \textbf{not exposed} by the current
    \filename{\_normalize\_usage} schema
    (\filename{src/llm\_client.py:42-63} records input / output /
    total only). The premium row is therefore marked
    \textbf{unavailable in v3}; a schema extension
    (\texttt{cache\_creation\_input\_tokens} /
    \texttt{cache\_read\_input\_tokens} / Gemini
    \texttt{thinking}-block) is flagged in \S\ref{sec:conclusion}.
\end{itemize}

A practitioner sizing budget against the public per-token rate
(input + output only) for a Kimi-/MiniMax-style vendor under-budgets
by exactly $r/(1-r)$ (equivalently, the ratio of hidden reasoning
tokens to visible completion tokens). This is the labelled sharp
observation the
introduction promised, surfaced as the single number most likely to
be cited in production deployment.

\subsection{Operating-point recommendations}
\label{sec:results-ops}

With every checkpoint hitting 100\% success on LC\_VCO\_20G the
recommendation collapses to a cost-only ordering: at $\leq$\$0.005/run
\textbf{DeepSeek V4-flash} is recommended with GPT-5.4-mini the
US fallback; in the \$0.02--0.03/run band the four China-frontier
non-flash checkpoints tie with Gemini 2.x and Haiku 4.5 (vendor
choice reduces to operational maturity, \S\ref{sec:disc-ops}); and
Opus 4.7 at a 233$\times$ premium adds no measurable quality
benefit on this workload.
